\section{Introduction} \label{sec:intro}

Learning transferrable spatiotemporal representations is a critical research area in computer vision, holding diverse applications across domains such as video search \cite{gabeur2020multi}, game control \cite{bruce2024genie}, robotic learning \cite{palme}, self-driving \cite{zablocki2022explainability}, and scientific studies~\cite{team2023gemini}. Recently, the advancement of Large Language Models (LLMs)~\cite{gpt3,gpt4,llama1,llama2} and their multimodal variations (MLLMs) \cite{gpt4v,mmgpt,llava,team2023gemini} have had a profound impact on vision research and other disciplines. Embedding videos effectively into these large models and harnessing their capabilities to enhance video understanding performance has emerged as pivotal tasks~\cite{li2023videochat,videochatgpt}.

Previous research has identified several effective learning schemes for video representations, including reconstructing videos with masked inputs~\cite{he2022masked,videomae,wang2023videomae,st_mae}, aligning videos with languages \cite{cpd,videoclip,videococa,umt}, and predicting the next token using videos \cite{flamingo,sun2023emu1,li2023mvbench}. These approaches have turned out to be complementary and can be unified through a progressive training scheme. Notably, methods such as InternVideo~\cite{wang2022internvideo}, UMT~\cite{umt}, and VideoPrism~\cite{videoprism} have utilized a two-stage training approach involving masked reconstruction and multimodal contrastive learning, leading to improved performance in downstream tasks. Following this line, we aim to further extend this progressive training scheme by incorporating video-based next token prediction and scaling the entire training process, including models and data, to build a new family of video foundation models. 

The proposed video foundation model, coined as {\modelname}, is built through a progressive training scheme. The learning involves three stages: (1) capturing spatiotemporal structure via unmasked reconstruction, (2) aligning with semantics from other modalities, and (3) enhancing its open-ended dialogue power through next token prediction.
In the initial stage, the model learns to reconstruct the unmasked video tokens, allowing the video encoder to develop basic spatiotemporal perception capability. To estimate the existing tokens, vision encoders (InternViT \cite{chen2023internvl} and VideoMAE-g \cite{wang2023videomae}) trained differently are employed as proxies.
In the next stage of crossmodal learning, the architecture is expanded to include audio and text encoders. This not only improves the alignment between videos and text but also endows {\modelname} the ability to handle video-audio tasks. By incorporating these additional modalities, the model's understanding of videos is enriched and aligned with their semantics.
Finally, in the next-token prediction stage, a video-centric dialogue system and the corresponding instruction-finetuning dataset are built to further tune the {\modelname}. By connecting {\modelname} to LLMs, the video encoder is further updated through next-token prediction training, enhancing its ability for open-ended tasks such as VQA and video description. 

For the training of {\modelname}, we emphasize the spatiotemporal consistency and labeling quality in the data. We build a large-scale multimodal video-centric dataset consisting of 402M data entries, which includes 2M videos, 50M video-text pairs (from WebVid~\cite{webvid} and InternVid~\cite{wang2023internvid}), 50M video-audio-speech-text pairs (InternVid2), and 300M image-text pairs.
Specifically, for InternVid2, we segment videos into clips semantically and focus on recalibrating the clip descriptions using three modalities: audio, video, and speech. We first generate captions for these three modalities separately. Then individual captions are fused together to create a more comprehensive description, which will improve the model's ability to comprehend and interpret the video accurately.

We evaluate {\modelname} across a wide range of video-related tasks. These tasks span from basic spatiotemporal perception, such as action recognition, to high-level reasoning tasks, such as long video or procedure-aware question-answering (QA), as given in Fig. \ref{fig:teaser}. The results (in Sec. \ref{sec:exp}) demonstrate that {\modelname} achieves the state-of-the-art performance on multiple tasks and is able to analyze and reason over sequences of actions. This top performance signifies its capability to effectively capture and understand video content. These empirical findings validate that {\modelname} could serves as a general video encoder for future exploration in video understanding. In summary, our contributions to video understanding are as follows.
\begin{itemize}[leftmargin=*]
\item[$\bullet$] This paper introduces {\modelname}, a competitive family of video foundation models that leverages masked reconstruction, crossmodal contrastive learning, and next token prediction to make model perceptive, semantic, and capable of reasoning in video understanding. 
\item[$\bullet$] {\modelname} achieves the state-of-the-art performance for more than 60 video / audio tasks. Our model demonstrates superior performance in video-related dialogue and long video understanding, highlighting its potential in modeling high-level world knowledge.
\item[$\bullet$] We provide an enhanced dataset to train {\modelname}. This includes the validation and incorporation of audio data during training, as well as the improved captioning method. These improvements result in significant enhancements in model performance and generalization ability.
\end{itemize}